% !PRINTED 185
\weberpage{185}
%<<CONT
erh\"alt man sechs Relationen zwischen den neun Coëfficienten
von $A$.

Diese Relationen lauten, wenn

\weq{3}{A = \begin{pmatrix} a_1, & a_2, & a_3\\ b_1, & b_2, & b_3\\ c_1, & c_2, & c_3 \end{pmatrix}}

angenommen wird:

\weq{4}{\begin{array}{r@{}c@{}l@{\qquad}r@{}c@{}l}
a_1^2+b_1^2+c_1^2 &{}={}& 1, & a_2a_3+b_2b_3+c_2c_3 &{}={}& 0\\
a_2^2+b_2^2+c_2^2 &{}={}& 1, & a_3a_1+b_3b_1+c_3c_1 &{}={}& 0\\
a_3^2+b_3^2+c_3^2 &{}={}& 1, & a_1a_2+b_1b_2+c_1c_2 &{}={}& 0,
\end{array}}

und sind aus der analytischen Geometrie wohl bekannt. F\"ur das
Quadrat der Substitutionsdeterminante $|A|$ ergiebt sich aus diesen
Relationen nach der Multiplicationsregel der Determinanten der
Werth 1, und folglich hat $|A|$ den Werth $\pm 1$.

Aus den Formeln (4) ergiebt sich, dass die \gesperrt{inverse Substitution} zu $A$

\[
A^{-1} = \begin{pmatrix} a_1, & b_1, & c_1\\ a_2, & b_2, & c_2\\ a_3, & b_3, & c_3 \end{pmatrix}
\]

mit der transponirten identisch ist, und diese Eigenschaft k\"onnte
auch als Definition der orthogonalen Substitutionen dienen.

Die Gesammtheit der orthogonalen Substitutionen bildet
eine Gruppe. Darunter ist eine engere Gruppe enthalten, die
durch den Werth

\weq{5}{|A| = + 1}

ausgezeichnet ist, die wir als die Gruppe der \gesperrt{eigentlichen
orthogonalen Substitutionen} bezeichnen wollen.

Man kann die lineare Substitution (1) als den Uebergang
von einem rechtwinkeligen Coordinatensysteme zu einem zweiten
mit demselben Anfangspunkte deuten, wenn man $x_1, x_2, x_3$ und
$y_1, y_2, y_3$ als rechtwinkelige Coordinaten eines und desselben
(ver\"anderlichen) Punktes in dem ersten und dem zweiten Coordinatensysteme ansieht. Durch $A$ ist die gegenseitige Lage der
beiden Coordinatensysteme bestimmt und umgekehrt. Besteht
die Bedingung (5), wie wir jetzt voraussetzen wollen, so kann
das erste Coordinatensystem mit dem zweiten zur Deckung gebracht werden durch Drehung um eine feste Axe mit einem
bestimmten Drehungswinkel.
